\documentclass[10pt,a4paper]{article}

\usepackage[utf8]{inputenc}
\usepackage[dutch]{babel}
\usepackage{fancyhdr}
\usepackage{geometry}
\usepackage{graphicx}
\usepackage{tabularx}
\usepackage{wallpaper}
\usepackage{listings}
\usepackage{amsmath}
\usepackage{amsthm}
\usepackage{amssymb}

\usepackage{xcolor, colortbl}
\definecolor{ugentblue}{HTML}{164A7C}
\definecolor{gray}{HTML}{AAAAAA}

% Tables
\def\arraystretch{1.35}
\renewcommand{\tabularxcolumn}[1]{>{\small}m{#1}}
\newcommand{\hcell}[1]{
	\cellcolor{ugentblue}\color{white}\textbf{#1}
}

\makeatletter
\renewcommand\thesubsection{\@arabic\c@section.\@arabic\c@subsection}
\makeatother{}

\usepackage[hypertexnames=false]{hyperref}
\usepackage[numbered, depth=3]{bookmark}

\renewcommand{\headrulewidth}{0pt}
\pagestyle{fancy}
\fancyhf{}

\newtheorem{theorem}{Theorem}

\interfootnotelinepenalty=0

\begin{document}
	\begin{titlepage}
		%% Footer
		\thispagestyle{fancy}
		\fancyhf{}
		
		%% Page
		\hfill
		\begin{minipage}[t][0.9\textheight]{0.8\textwidth}
			\noindent
			\includegraphics[width=55px]{ugent-blue.png} \\[-1em]
			\color{ugentblue}
			\makebox[0pt][l]{\rule{1.3\textwidth}{1pt}}
			\par
			\noindent
			\textbf{\textsf{Formele Logica}} \textcolor{gray}{\textsf{Academiejaar 2014-2015}}
			\vfill
				\noindent
			{\huge \textsf{Project Alloy}}
			\vskip\baselineskip
			\noindent
			\textsf{\textbf{Jasper D'haene}}
			\textsf{\textbf{Florian Dejonckheere}}
		\end{minipage}
	\end{titlepage}

	%%% PAGE STYLE %%%
	\nopagecolor
	\renewcommand{\footrulewidth}{0.4pt}
	\headheight 45pt
	\ULCornerWallPaper{1}{header.png}
	\fancyfoot[C]{\thepage}

	%%% DOCUMENT %%%
	\section{Ramsey-getallen}
		Ramsey's theorem dicteert dat in een voldoende grote set waarvan de bogen gekleurd zijn met een willekeurig aantal kleuren, monochromatisch gekleurde subsets te vinden zijn. Stel $m, n \in \mathbb{R^+}$ en twee kleuren $k_i, i \in [0, 1]$, dan definieert Ramsey's theorem $R(m, n, k)$ de ondergrens voor de complete graaf die een subset over tenminste $m$ toppen met kleur $k_0$, of een subset over tenminste $n$ toppen met kleur $k_1$ bevat.\\

		\noindent Informeel is dit probleem (voor $k = 2$) ook wel bekend als het \textit{party problem}: hoe groot is de minimale set van personen die uitgenodigd moeten worden voor een feestje, waarvoor geldt dat ofwel minstens $m$ personen elkaar (mutueel) kennen, ofwel $n$ personen elkaar (mutueel) niet kennen.

	\clearpage{}
	\section{Cyclische toren van Hanoi}
		De cyclische toren van Hanoi is een variant op de bekende combinatorische puzzel. Buiten de drie hoofdregels geldt de volgende regels ook:
		\begin{enumerate}
			\setcounter{enumi}{3}
			\item Schijven kunnen enkel cyclisch naar rechts opgeschoven worden. Voor drie palen $A, B, C$ is dit dus $A \rightarrow B \rightarrow C \rightarrow A$
		\end{enumerate}

		\noindent Als basis voor het algoritme werd de ge\"implementeerde versie van Ilya Shlyakhter beschouwd. Deze versie wordt bijgeleverd als voorbeeld bij de Alloy Analyzer, maar implementeert de niet-cyclische variant.

\end{document}
